% STATUS: rewritten 2026-11-05; every claim is cited.
\section{Semirings}\label{sec:other-systems:semirings}

The $\Gamma$-rings of the previous section still had additive negatives,
and a semiring drops them: proofs that use subtraction then break, and
the methods change \cite[p.~624]{AhmedDudek2021}. Ahmed and Dudek
describe the setting as follows
\cite[p.~625]{AhmedDudek2021}.

\begin{definition}[{\cite[p.~625]{AhmedDudek2021}}]\label{def:other:semiring}
A \emph{semiring} $(S, +, \cdot)$ is a nonempty set $S$ with two binary
operations $+$ and $\cdot$ such that $\cdot$ is left and right
distributive over $+$, $(S, +)$ is a semigroup with neutral element $0$,
and $(S, \cdot)$ is a semigroup with zero $0$, that is $0\,a = a\,0 = 0$
for all $a \in S$. If $(S, \cdot)$ is commutative, the semiring $S$ is
called \emph{commutative}.
\end{definition}

The next definition states the additively inverse semirings.

\begin{definition}[{\cite[p.~625]{AhmedDudek2021}}]\label{def:other:ma-semiring}
A semiring $S$ is \emph{additively inverse} if for every $a \in S$ there
is a unique $a' \in S$ such that $a + a' + a = a$ and $a' + a + a' = a'$.
An additively inverse semiring $S$ with commutative addition such that
$a + a'$ lies in the center $Z(S)$ for every $a \in S$ is called an
\emph{M A-semiring}.
\end{definition}

The next definition states the derivations on an additively inverse semiring.

\begin{definition}[{\cite[p.~626]{AhmedDudek2021}}]\label{def:other:semiring-derivations}
Let $S$ be an additively inverse semiring. An additive mapping
$d \colon S \to S$ is a \emph{derivation} if
$d(xy) = d(x)\,y + x\,d(y)$ for all $x, y \in S$, and a \emph{Jordan
derivation} if $d(x^{2}) = d(x)\,x + x\,d(x)$ for all $x \in S$. It is a
\emph{left derivation} if $f(xy) = x\,f(y) + y\,f(x)$, and a \emph{left
Jordan derivation} if
\[
f(x^{2}) + 2\,x'\,f(x) = 0 \qquad (x \in S).
\]
\end{definition}

Ahmed and Dudek classify the prime M A-semirings on which a nonzero
left Jordan derivation exists: they are commutative
\cite[pp.~624--631]{AhmedDudek2021}. Their main theorem is the analogue of
the result of Bre\v{s}ar and Vukman on prime rings.

\begin{theorem}[{\cite[Theorem~4.3, p.~631]{AhmedDudek2021}}]\label{thm:other:ahmed-dudek}
If on a prime, $2$- and $3$-torsion free M A-semiring $S$ there exists a
non-zero left Jordan derivation, then $S$ is commutative.
\end{theorem}

The proof replaces subtraction by the additive-inverse map $a \mapsto a'$
and keeps the torsion conditions to make the commutator work
\cite[pp.~631--632]{AhmedDudek2021}. The next section keeps the order but
drops the product and moves to derivations on lattices
\cite[Definition~3.1, p.~309]{XinLiLu2008}.

Vladeva studies the Jordan derivations of the finite endomorphism semiring,
where the additive structure is a semigroup rather than a group
\cite[Abstract, p.~1]{Vladeva2017Endomorphism}.

Dai and DeMarr show that the range of a positive derivation of a Dedekind $\sigma$-complete partially ordered linear algebra with an order unit consists of generalized nilpotents, and that with further assumptions the algebra has a property like that of the algebra of upper triangular matrices~\cite{WlPositiveDerivationsOnPartially1978}.

Ali, Aslam, Qureshi, Khan, Rehman and Farid study commutativity of MA-semirings with involution through generalized derivations~\cite{WlCommutativityOfMaSemiringsWith2020}.

Ali and Aslam study generalized derivations of MA-semirings with involution that satisfy stated identities, using Hermitian and skew-Hermitian elements, and give conditions under which the semiring is commutative~\cite{WlOnCertainIdentitiesOfGeneraliz2026}.
